%% ma: L10-B15-0050
%% lop: 10
%% chuong: 6
%% bai: 15
%% dang: 10.15.5
%% loai: TL
%% muc: VD
%% dap_an: "-"
%% nguon: "Ảnh giáo viên gửi 21/11/2026: Chuyên đề VI – Toán 10 – Chương VI (Hàm số – đồ thị và ứng dụng), tr. 199–203, Dạng 5 Câu 3"
%% kien_thuc: "Lớp 10 Bài 15 (hàm số đồng biến, nghịch biến)"
%% trang_thai: chinh-thuc
%% ghi_chu: "Dạng do giáo viên cung cấp (21/11/2026)."
\begin{ex}
Tìm tất cả các giá trị thực của tham số $m$ để hàm số $f(x)=-x^2+(m-1)x+2$ nghịch biến trên khoảng $(1;2)$.
\loigiai{Với $x_1\ne x_2$: $\dfrac{f(x_1)-f(x_2)}{x_1-x_2}=-(x_1+x_2)+m-1$.

Cần $m-1<x_1+x_2$ với mọi $x_1\ne x_2$ thuộc $(1;2)$. Khi đó $x_1+x_2$ nhận mọi giá trị trong $(2;4)$, nên điều kiện là $m-1\le2\Leftrightarrow m\le3$.}
\end{ex}
